\documentclass[11pt,a4paper]{article}
\usepackage[utf8]{inputenc}
\usepackage[T1]{fontenc}
\usepackage[brazil]{babel}
\usepackage{geometry}
\usepackage{enumitem}
\usepackage{xcolor}
\usepackage{titlesec}
\usepackage{fancyhdr}
\usepackage{listings}
\usepackage{tcolorbox}

\geometry{margin=2.5cm}

% Cores
\definecolor{headerblue}{RGB}{0,102,204}
\definecolor{sectiongray}{RGB}{100,100,100}
\definecolor{codebg}{RGB}{245,245,245}

% Estilo de seções
\titleformat{\section}{\Large\bfseries\color{headerblue}}{}{0em}{}
\titleformat{\subsection}{\large\bfseries\color{sectiongray}}{}{0em}{}

% Cabeçalho e rodapé
\pagestyle{fancy}
\fancyhf{}
\fancyhead[L]{\textbf{SIVIRA - Factory}}
\fancyhead[R]{\thepage}
\renewcommand{\headrulewidth}{0.5pt}

% Configuração de código
\lstset{
    basicstyle=\ttfamily\small,
    backgroundcolor=\color{codebg},
    breaklines=true,
    frame=single,
    numbers=left,
    numberstyle=\tiny\color{gray}
}

\begin{document}

% Capa
\begin{titlepage}
    \centering
    \vspace*{3cm}

    {\Huge\bfseries Sistema Inteligente para Visualização\par}
    {\Huge\bfseries e Replanejamento de Atividades\par}
    \vspace{0.5cm}
    {\Large\bfseries (SIVIRA)\par}

    \vspace{2cm}
    {\LARGE\bfseries Documentação do Módulo Factory\par}

    \vspace{2cm}
    {\large Módulo: \texttt{factory/}\par}

    \vspace{1cm}
    {\large Padrão de Design: Factory Pattern + Singleton\par}

    \vfill

    {\large \today\par}
\end{titlepage}

\tableofcontents
\newpage

% Introdução
\section{Introdução}

O módulo \texttt{factory/} implementa o padrão de design Factory Pattern para criação centralizada de objetos no sistema SIVIRA. Este módulo é responsável por:

\begin{itemize}
    \item Criar instâncias de equipamentos com configurações padronizadas
    \item Gerenciar funcionários a partir de arquivos de configuração JSON
    \item Garantir consistência na criação de objetos através do padrão Singleton
    \item Facilitar a manutenção e evolução do sistema
\end{itemize}

\subsection{Arquivos do Módulo}

\begin{itemize}
    \item \texttt{fabrica\_equipamentos.py} (782 linhas): Factory para criação de equipamentos
    \item \texttt{fabrica\_funcionarios.py} (342 linhas): Factory para gerenciamento de funcionários
    \item \texttt{\_\_init\_\_.py}: Módulo de inicialização
\end{itemize}

\newpage

% Fábrica de Equipamentos
\section{Fábrica de Equipamentos}

\subsection{Visão Geral}

A classe \texttt{FabricaEquipamentos} é responsável pela criação padronizada de todos os equipamentos do sistema. Utiliza métodos estáticos para instanciar equipamentos com configurações específicas.

\subsection{Estrutura da Classe}

\begin{tcolorbox}[colback=codebg,colframe=headerblue,title=Classe FabricaEquipamentos]
\begin{lstlisting}[language=Python]
class FabricaEquipamentos:
    """
    Fabrica responsavel pela criacao dos equipamentos.
    """
\end{lstlisting}
\end{tcolorbox}

\subsection{Equipamentos Criados}

A fábrica cria 17 tipos diferentes de equipamentos, organizados por categoria:

\subsubsection{Refrigeração e Congelamento}

\begin{itemize}[leftmargin=2cm]
    \item \textbf{criar\_camara\_refrigerada\_1()}: Câmara com temperatura 0-4°C, capacidade 200 caixas
    \item \textbf{criar\_camara\_refrigerada\_2()}: Câmara com temperatura -18 a -7°C, capacidade 200 caixas
    \item \textbf{criar\_freezer\_1()}: Freezer com temperatura -25 a -10°C
\end{itemize}

\subsubsection{Cocção}

\begin{itemize}[leftmargin=2cm]
    \item \textbf{criar\_fogao\_1()}: Fogão de cozinha, 6 bocas, alta pressão
    \item \textbf{criar\_fogao\_2()}: Fogão de confeitaria, 4 bocas, baixa pressão
    \item \textbf{criar\_forno\_1()}: Forno turbo/lastro, 4 níveis, temperatura até 300°C
    \item \textbf{criar\_forno\_2()}: Forno combinado, 5 níveis, temperatura até 250°C
    \item \textbf{criar\_fritadeira\_1()}: Fritadeira de 15L com 2 frações
\end{itemize}

\subsubsection{Mistura e Preparação}

\begin{itemize}[leftmargin=2cm]
    \item \textbf{criar\_batedeira\_industrial\_1()}: Capacidade 5-40kg, velocidade 100-500 RPM
    \item \textbf{criar\_batedeira\_planetaria\_1()}: Capacidade 1-7kg, 3 velocidades
    \item \textbf{criar\_masseira\_1()}: Capacidade 15-50kg, mistura rápida/semi-rápida
    \item \textbf{criar\_hot\_mix\_1()}: Capacidade 5-30kg com cocção integrada
\end{itemize}

\subsubsection{Modelagem e Divisão}

\begin{itemize}[leftmargin=2cm]
    \item \textbf{criar\_divisora\_de\_massas\_1()}: Capacidade 30-3000g, 30 divisões
    \item \textbf{criar\_modeladora\_de\_paes\_1()}: Modelagem de pães, capacidade 30-300g
    \item \textbf{criar\_modeladora\_de\_salgados\_1()}: Modelagem de salgados, capacidade 20-200g
\end{itemize}

\subsubsection{Fermentação e Armazenamento}

\begin{itemize}[leftmargin=2cm]
    \item \textbf{criar\_armario\_fermentador\_1()}: Temperatura 30-40°C, 9 níveis
    \item \textbf{criar\_armario\_esqueleto\_1()}: Armazenamento com 32 níveis
\end{itemize}

\subsubsection{Pesagem e Embalagem}

\begin{itemize}[leftmargin=2cm]
    \item \textbf{criar\_balanca\_digital\_1()}: Pesagem 1-30000g, precisão 1g
    \item \textbf{criar\_embaladora\_1()}: Embalagem a vácuo, capacidade 100-5000g
\end{itemize}

\subsubsection{Bancadas}

\begin{itemize}[leftmargin=2cm]
    \item \textbf{criar\_bancada\_1()}: Bancada com 4 frações divisíveis
\end{itemize}

\subsection{Características Técnicas}

\subsubsection{Capacidades}

Cada equipamento possui capacidades mínimas e máximas bem definidas:

\begin{itemize}
    \item \textbf{Misturadoras}: 1kg a 50kg dependendo do tipo
    \item \textbf{Fornos}: 4 a 5 níveis com temperaturas de 150°C a 300°C
    \item \textbf{Câmaras}: 200 caixas com controle de temperatura
    \item \textbf{Fritadeiras}: 15L com sistema de frações
\end{itemize}

\subsubsection{Setores}

Os equipamentos são distribuídos entre os setores:

\begin{itemize}
    \item \textbf{ALMOXARIFADO}: Câmaras refrigeradas, freezer
    \item \textbf{COZINHA}: Fogão 1, fritadeira
    \item \textbf{CONFEITARIA}: Fogão 2, batedeiras, hot mix
    \item \textbf{PANIFICACAO}: Fornos, masseira, modeladoras, fermentador
\end{itemize}

\subsection{Correções Implementadas}

Durante o desenvolvimento, foram realizadas as seguintes correções:

\begin{itemize}
    \item Capacidade de câmaras refrigeradas: 20 → 200 caixas
    \item Faixa de temperatura: ajustada para valores operacionais reais
    \item Número de operadores de fogões: ajustado para capacidade real
    \item Capacidades mínimas de bocas: reduzidas para permitir batches pequenos (2000g → 500g)
\end{itemize}

\newpage

% Fábrica de Funcionários
\section{Fábrica de Funcionários}

\subsection{Visão Geral}

A classe \texttt{FabricaFuncionarios} implementa o padrão Factory combinado com Singleton para gerenciar funcionários. Diferente da fábrica de equipamentos, esta carrega dados de um arquivo JSON, permitindo fácil manutenção sem modificar código.

\subsection{Padrão Singleton}

\begin{tcolorbox}[colback=codebg,colframe=headerblue,title=Implementação do Singleton]
\begin{lstlisting}[language=Python]
class FabricaFuncionarios:
    """
    Factory para criacao de funcionarios a partir de arquivo JSON.

    Singleton: Garante que apenas uma instancia da fabrica existe,
    mantendo os mesmos objetos Funcionario em memoria durante toda
    a execucao.
    """

    _instance = None
    _funcionarios_carregados = False

    def __new__(cls, caminho_json: str = ...):
        if cls._instance is None:
            cls._instance = super(FabricaFuncionarios, cls).__new__(cls)
        return cls._instance
\end{lstlisting}
\end{tcolorbox}

\subsection{Características do Singleton}

\begin{itemize}
    \item \textbf{Instância única}: Garante apenas um objeto FabricaFuncionarios em memória
    \item \textbf{Estado preservado}: Funcionários mantêm ocupações entre múltiplas chamadas
    \item \textbf{Evita recarregamento}: Dados JSON carregados apenas uma vez
    \item \textbf{Consistência}: Mesmos objetos Funcionario em toda aplicação
\end{itemize}

\subsection{Métodos Principais}

\subsubsection{carregar\_funcionarios()}

Carrega funcionários do arquivo JSON. Se já foram carregados anteriormente (singleton), retorna a lista existente preservando ocupações.

\begin{tcolorbox}[colback=codebg,colframe=headerblue,title=Método carregar\_funcionarios]
\begin{lstlisting}[language=Python]
def carregar_funcionarios(self) -> List[Funcionario]:
    """
    Carrega funcionarios do arquivo JSON.

    Singleton: Se ja foi carregado antes, retorna a mesma lista
    de funcionarios (com ocupacoes preservadas).
    """
    if FabricaFuncionarios._funcionarios_carregados:
        return self.funcionarios_disponiveis

    # Carrega do JSON...
    FabricaFuncionarios._funcionarios_carregados = True
    return self.funcionarios_disponiveis
\end{lstlisting}
\end{tcolorbox}

\subsubsection{obter\_funcionario\_por\_id(id)}

Busca funcionário específico por ID.

\subsubsection{obter\_funcionarios\_por\_setor(setor)}

Retorna lista de funcionários que trabalham em um setor específico.

\subsubsection{obter\_funcionarios\_por\_profissao(profissao)}

Retorna lista de funcionários com uma profissão específica.

\subsection{Estrutura do JSON}

O arquivo \texttt{data/funcionarios/funcionarios.json} segue esta estrutura:

\begin{tcolorbox}[colback=codebg,colframe=headerblue,title=Estrutura JSON]
\begin{lstlisting}[language=JSON]
{
  "funcionarios": [
    {
      "id": 1,
      "nome": "Funcionario 1",
      "setores": ["PANIFICACAO", "CONFEITARIA"],
      "tipos_profissionais": ["PADEIRO"],
      "regras_folga": [
        {
          "tipo": "DIA_FIXO_SEMANA",
          "dia_semana": "SEXTA"
        }
      ],
      "ch_semanal": 44,
      "horario_inicio": "08:00",
      "horario_final": "18:00",
      "horario_intervalo": {
        "inicio": "11:00",
        "duracao_minutos": 60
      },
      "fip": 2.0
    }
  ]
}
\end{lstlisting}
\end{tcolorbox}

\subsection{Campos do Funcionário}

\begin{itemize}[leftmargin=2cm]
    \item \textbf{id}: Identificador único
    \item \textbf{nome}: Nome do funcionário
    \item \textbf{setores}: Lista de setores onde pode trabalhar
    \item \textbf{tipos\_profissionais}: Lista de profissões/habilidades
    \item \textbf{regras\_folga}: Padrões de folga (dia fixo, n-ésimo dia do mês, etc.)
    \item \textbf{ch\_semanal}: Carga horária semanal
    \item \textbf{horario\_inicio}: Horário de entrada
    \item \textbf{horario\_final}: Horário de saída
    \item \textbf{horario\_intervalo}: Início e duração do intervalo
    \item \textbf{fip}: Fator de Importância do Profissional
\end{itemize}

\subsection{Tipos de Folga Suportados}

\begin{itemize}[leftmargin=2cm]
    \item \textbf{DIA\_FIXO\_SEMANA}: Folga em dia fixo da semana (ex: toda segunda-feira)
    \item \textbf{DIA\_FIXO\_MES}: Folga em dia fixo do mês (ex: todo dia 15)
    \item \textbf{N\_DIA\_SEMANA\_DO\_MES}: Folga no n-ésimo dia da semana do mês (ex: segunda segunda-feira)
\end{itemize}

\subsection{Funcionários Cadastrados}

O sistema possui 9 funcionários cadastrados:

\begin{enumerate}
    \item \textbf{Funcionário 1}: Padeiro (Panificação/Confeitaria) - FIP 2.0
    \item \textbf{Funcionário 2}: Auxiliar de Padeiro (Panificação) - FIP 1.0
    \item \textbf{Funcionário 3}: Auxiliar de Padeiro/Confeiteiro (Multi-setores) - FIP 3.0
    \item \textbf{Funcionário 4}: Confeiteiro (Confeitaria) - FIP 2.0
    \item \textbf{Funcionário 5}: Auxiliar de Confeiteiro (Confeitaria) - FIP 2.0
    \item \textbf{Funcionário 6}: Auxiliar Multi-função (Confeitaria/Panificação) - FIP 1.0
    \item \textbf{Funcionário 7}: Cozinheiro (Cozinha) - FIP 1.0
    \item \textbf{Funcionário 8}: Almoxarife (Almoxarifado) - FIP 1.0
    \item \textbf{Funcionário 9}: Almoxarife (Almoxarifado) - FIP 1.0
\end{enumerate}

\subsection{Compatibilidade com Código Legado}

Para manter compatibilidade com código existente, a fábrica cria variáveis globais individuais:

\begin{lstlisting}[language=Python]
funcionario_1 = funcionarios_disponiveis[0]
funcionario_2 = funcionarios_disponiveis[1]
# ... etc
\end{lstlisting}

Isso permite que código antigo continue funcionando enquanto novo código pode usar a API da fábrica.

\newpage

\section{Integração com o Sistema}

\subsection{Uso da Fábrica de Equipamentos}

\begin{lstlisting}[language=Python]
from factory.fabrica_equipamentos import FabricaEquipamentos

# Criar equipamento especifico
forno = FabricaEquipamentos.criar_forno_1()

# Acessar propriedades
print(forno.nome)
print(forno.capacidade_niveis_max)
\end{lstlisting}

\subsection{Uso da Fábrica de Funcionários}

\begin{lstlisting}[language=Python]
from factory.fabrica_funcionarios import FabricaFuncionarios

# Criar instancia (singleton)
fabrica = FabricaFuncionarios()

# Carregar funcionarios
funcionarios = fabrica.carregar_funcionarios()

# Buscar por setor
panificacao = fabrica.obter_funcionarios_por_setor(
    TipoSetor.PANIFICACAO
)

# Buscar por profissao
padeiros = fabrica.obter_funcionarios_por_profissao(
    TipoProfissional.PADEIRO
)
\end{lstlisting}

\subsection{Benefícios do Padrão Singleton}

\begin{itemize}
    \item \textbf{Memória eficiente}: Evita múltiplas instâncias desnecessárias
    \item \textbf{Estado consistente}: Ocupações de funcionários persistem durante execução
    \item \textbf{Performance}: JSON carregado apenas uma vez
    \item \textbf{Simplicidade}: Interface única e centralizada
\end{itemize}

\newpage

\section{Estatísticas}

\subsection{Resumo do Módulo}

\begin{tabular}{lr}
\textbf{Métrica} & \textbf{Valor} \\
\hline
Total de linhas de código & 1.124 \\
Arquivos Python & 2 \\
Equipamentos criáveis & 17 \\
Funcionários cadastrados & 9 \\
Setores suportados & 4 \\
Tipos de folga & 3 \\
\hline
\end{tabular}

\subsection{Distribuição por Arquivo}

\begin{tabular}{lrr}
\textbf{Arquivo} & \textbf{Linhas} & \textbf{\%} \\
\hline
fabrica\_equipamentos.py & 782 & 69.6\% \\
fabrica\_funcionarios.py & 342 & 30.4\% \\
\hline
\textbf{Total} & \textbf{1.124} & \textbf{100\%} \\
\end{tabular}

\subsection{Equipamentos por Setor}

\begin{tabular}{lr}
\textbf{Setor} & \textbf{Equipamentos} \\
\hline
Panificação & 7 \\
Confeitaria & 5 \\
Cozinha & 2 \\
Almoxarifado & 3 \\
\hline
\textbf{Total} & \textbf{17} \\
\end{tabular}

\newpage

\section{Padrões de Design}

\subsection{Factory Pattern}

Ambas as fábricas implementam o padrão Factory, que:

\begin{itemize}
    \item Centraliza a lógica de criação de objetos
    \item Encapsula detalhes de implementação
    \item Facilita manutenção e evolução
    \item Permite substituição de implementações
\end{itemize}

\subsection{Singleton Pattern}

A FabricaFuncionarios adiciona o padrão Singleton para:

\begin{itemize}
    \item Garantir instância única
    \item Preservar estado durante execução
    \item Evitar recarregamento desnecessário
    \item Manter consistência de dados
\end{itemize}

\subsection{Separation of Concerns}

O módulo separa claramente:

\begin{itemize}
    \item \textbf{Configuração}: Dados JSON (funcionários)
    \item \textbf{Criação}: Métodos factory (equipamentos)
    \item \textbf{Modelo}: Classes de domínio (models/)
    \item \textbf{Negócio}: Lógica de alocação (services/)
\end{itemize}

\newpage

\section{Conclusão}

O módulo \texttt{factory/} implementa uma arquitetura robusta e extensível para criação de objetos no sistema SIVIRA. A combinação dos padrões Factory e Singleton garante:

\begin{itemize}
    \item \textbf{Consistência}: Objetos criados de forma padronizada
    \item \textbf{Manutenibilidade}: Configurações centralizadas e fáceis de modificar
    \item \textbf{Performance}: Singleton evita overhead de múltiplas instâncias
    \item \textbf{Flexibilidade}: JSON permite mudanças sem recompilar
    \item \textbf{Escalabilidade}: Fácil adicionar novos equipamentos ou funcionários
\end{itemize}

\subsection{Próximos Passos}

Possíveis melhorias futuras:

\begin{enumerate}
    \item Migrar equipamentos para configuração JSON
    \item Adicionar validação de schemas JSON
    \item Implementar cache persistente entre execuções
    \item Adicionar suporte a múltiplos perfis de configuração
    \item Criar interface gráfica para edição de funcionários
\end{enumerate}

\end{document}
